% =====================================================================
%  MARCO REFERENCIAL
%
%  Reescrito el 2026-10-15. Cambios respecto a la versión anterior:
%   - Los antecedentes suman la línea de robots que guían personas
%     (Burgard 1999, Guerreiro 2019) y cierran con una tabla del aporte.
%   - El marco teórico se apoya en la literatura de cada decisión: taxonomías
%     de asignación (Gerkey 2004, Korsah 2013), cinemática y planificación
%     (Siegwart, Reeds-Shepp, Dolgov, Smac), seguimiento (Coulter, RPP),
%     localización (Fox, Thrun), odometría láser (Jaimez), degeneración
%     (Zhang), fusión (Moore), Gazebo, rosbridge, Wilson, Brown y Popper.
%   - Las referencias internas a secciones van con \label y \ref.
%   - «Hueco» y otras palabras coloquiales fuera, por indicación del director.
%
%  REGISTRO: lectura de pregrado. Frases cortas, voz activa, una idea por
%  frase. Los términos van en negrita solo en el marco conceptual, donde
%  sirven para barrer la sección con la vista (pedido del director).
% =====================================================================
\chapter{Marco referencial}
\label{cap:marco}

Este capítulo sitúa el trabajo frente a lo que ya existe y fija el vocabulario y
la teoría con que se lee el resto del documento. Los antecedentes revisan cinco
líneas de trabajo y terminan con lo que este proyecto aporta frente a ellas. El
marco conceptual define los términos técnicos. El marco teórico expone el
fundamento de cada decisión de diseño y la literatura que lo respalda. El marco
legal y el ambiental sitúan el prototipo frente a la normativa que le aplica y
frente a su impacto.

% =====================================================================
\section{Antecedentes}
\label{sec:antecedentes}
% =====================================================================
% Viene del capitulo 5 del anteproyecto (Revision del estado del arte), con
% citas de autor: se nombra a quien hizo el trabajo en lugar de dejar solo el
% numero entre corchetes. La plantilla pide «desarrollo de los antecedentes» en
% la introduccion, y el marco referencial es el capitulo donde caben.
% =====================================================================

La revisión se ordena en cinco líneas: robots que guían a personas, navegación
autónoma en interiores, coordinación entre varios robots, movilidad entre pisos
y distribución del cómputo. Cada una resuelve bien una parte del problema. La
última subsección reúne lo que ninguna resuelve, que es lo que este trabajo
aborda.

\subsection{Robots que guían a personas}

La idea de un robot que acompaña a una persona dentro de un edificio público
tiene más de veinticinco años. \textcite{burgard1999experiences} describen
RHINO, un robot guía que operó de forma autónoma durante seis días en un museo en
1997. Su arquitectura integraba localización, construcción de mapas, evasión de
obstáculos, planificación e interacción con los visitantes, además de una
interfaz web para seguirlo a distancia. De ese trabajo se toma una idea que sigue
vigente: guiar personas es un problema de integración tanto como de navegación,
porque ninguna de las piezas sirve si las demás fallan. Del mismo grupo salió poco
después la localización de Monte Carlo \autocite{fox1999monte}, que es la base
del algoritmo de localización que usan los vehículos de este proyecto.

\textcite{guerreiro2019cabot} trasladan la idea a la asistencia de personas
ciegas. Su robot, CaBot, tiene forma de maleta y transmite la dirección a seguir
por medio de un mango que vibra. Lo evaluaron con diez participantes ciegos, que
lo compararon con un perro guía o con un acompañante vidente. El trabajo muestra
dos exigencias que también aparecen aquí. El robot tiene que adaptarse a una
persona que camina a su lado. Y la interfaz tiene que comunicar el estado del
guiado sin exigir atención constante.

Compartir el espacio con personas plantea además un problema de seguridad.
\textcite{li2023bug} lo abordan modificando el algoritmo Bug con sensores de
proximidad, para reducir las colisiones entre un robot colaborativo y las
personas que trabajan a su alrededor. En este proyecto la respuesta es más
modesta y está en la operación: una persona acompaña a cada vehículo durante las
maniobras, y el usuario puede cancelar la misión en cualquier momento.

\textcite{jimenez2024arquitectura} integra al usuario dentro del propio sistema
de mapas, combinando representaciones semánticas con reconocimiento de personas
para localizar y seguir a quien acompaña, en una plataforma asistencial de varios
niveles. Es el antecedente más próximo en intención, porque trata a la persona
como parte del problema de navegación y no como un destinatario externo.

En los tres casos un mismo robot acompaña a la persona durante todo el recorrido.
Ese supuesto es el que este trabajo cambia.

\subsection{Navegación autónoma en interiores}

\textcite{sojan2025autonomous} integran percepción, localización y planificación
sobre ROS~2 para desplazamiento autónomo en interiores dinámicos. Su aporte es
consolidar el flujo completo (mapa, localización por filtro de partículas y
planificación) sobre una representación plana del espacio. Es la arquitectura
que este proyecto toma como punto de partida en cada vehículo.

El sistema de navegación que implementa ese flujo en ROS~2 es Nav2.
\textcite{macenski2020marathon} lo presentan como sucesor de la pila de
navegación de la primera generación de ROS. Organiza la tarea con un árbol de
comportamiento y gestiona el arranque de cada servidor con nodos de ciclo de
vida, que deben configurarse y activarse antes de funcionar. Lo validaron en un
campus universitario, con el robot navegando entre estudiantes durante una
prueba de larga duración. \textcite{macenski2023survey} recogen,
desde la experiencia de los propios mantenedores, los algoritmos de
planificación, control y localización que hoy ofrece el ecosistema, con pruebas
comparativas entre ellos. Las dos referencias sirvieron para elegir los
componentes que se describen en la sección~\ref{sec:teo-planificacion}.

\textcite{varghese2023multirobot} llevan el problema a un escenario cooperativo:
varios agentes construyen un mapa compartido y mantienen consistente su
localización. El énfasis se desplaza de la autonomía individual a la
construcción colectiva del entorno. Con el trabajo anterior comparte la técnica,
porque las dos propuestas usan representaciones bidimensionales y se apoyan en
ROS; la diferencia está en el alcance.

Existe una línea distinta que resuelve la navegación con aprendizaje en lugar de
con planificación explícita. \textcite{bojarski2016end} mostraron que una red
neuronal entrenada de extremo a extremo, solo con imágenes de cámara y el ángulo
de dirección del conductor, aprende a conducir sin que nadie programe el
seguimiento de carril. \textcite{zain2025obstacle} aplica una idea parecida a la
evasión de obstáculos, fusionando sensores con una red convolucional, y
\textcite{lopez2024enhancing} combinan aprendizaje profundo con lógica difusa
para obtener autonomía reactiva.

Este proyecto no sigue esa vía. Entrenar cualquiera de esos modelos exige un
conjunto de datos etiquetados del edificio concreto, que no existe y que habría
que levantar. Además, el aporte declarado aquí está en la coordinación de dos
agentes que se reparten un recorrido, y la percepción puede resolverse con
métodos establecidos.
Con planificación clásica cada decisión del vehículo se puede inspeccionar, y eso
pesa cuando hay que explicar un fallo. Una red entrenada podría dar mejor
desempeño reactivo, pero no permitiría explicar por qué tomó una decisión.

\subsection{Coordinación entre varios robots}

Las estrategias de coordinación de la literatura revisada se diferencian por el
nivel en el que cooperan. \textcite{nair2024collaborative} funda la cooperación
en el intercambio de información sensorial: los robots comparten lo que perciben
para mejorar el desempeño del conjunto. \textcite{roszkowska2023multilevel}
proponen en cambio una estructura jerárquica que supervisa el desplazamiento
simultáneo de varios agentes, con el fin de prevenir conflictos.
\textcite{rudrusamy2022broadcast} resuelven el problema desde el canal de
comunicación, con un esquema de control por difusión y seguimiento visual.

\textcite{xu2025modular} aportan una arquitectura modular sobre ROS para
asignación cooperativa de tareas, y \textcite{zhan2025kubernetes} estudian la
resiliencia del sistema cuando un nodo falla. El marco teórico de la asignación
de tareas lo fijan \textcite{gerkey2004taxonomy}, cuya taxonomía permite
clasificar cualquier esquema de reparto según cuántas tareas acepta un robot,
cuántos robots requiere una tarea y con cuánta anticipación se decide. La
sección~\ref{sec:teo-asignacion} usa esa taxonomía, y su extensión por
\textcite{korsah2013comprehensive}, para situar el esquema de este trabajo.

El antecedente más cercano a la propuesta de este proyecto es
\textcite{srivastava2025deliver}, que coordinan recogida y entrega entre varios
robots mediante planificación con relevos. Allí lo que se traspasa de un agente a
otro es un paquete; aquí es una persona. La diferencia tiene consecuencias de
diseño: un paquete puede esperar indefinidamente en el punto de transferencia, y
una persona no. Por eso el relevo de este trabajo incluye una confirmación
explícita del usuario cuando llega al otro piso, con un aviso a los 60~s y un
plazo máximo de 120~s.

Un antecedente de alcance comparable, por tratarse también de un proyecto de
grado de ingeniería, es \textcite{piguave2023multirobot}, que implementan un
sistema multi-robot de código abierto sobre ROS~2 para ambientes colaborativos.

Según la revisión del anteproyecto, todos estos trabajos operan en un solo nivel
físico, sin segmentación vertical.

\subsection{Movilidad entre pisos}

El cambio de nivel se ha abordado sobre todo desde la logística y la
exploración. \textcite{kim2024indoor} resuelven la planificación tridimensional
combinando un mapa por cada piso con un grafo topológico que los conecta, e
integran además la interacción con el ascensor. \textcite{song2024elevator} se
concentran en esto último: automatizar el uso del elevador mediante dispositivos
conectados instalados en el edificio.

\textcite{luo2025elevator} plantean la pregunta un paso antes (ascensor o
escalera) y la resuelven con aprendizaje por refuerzo jerárquico, de modo que el
robot elige el medio de transición según el tránsito que encuentra. En otra
dirección, \textcite{chen2025lite} estructuran el entorno con una descomposición
piso-escalera que extiende las técnicas bidimensionales a configuraciones de
varios niveles, y \textcite{gong2026stairway} incorporan conciencia de piso a la
navegación hacia objetos sin entrenamiento previo del escenario.

Las cuatro propuestas resuelven la movilidad entre pisos desde ángulos
complementarios: la infraestructura física, el modelado topológico y la decisión
adaptativa. Comparten un supuesto: el robot cruza de nivel, sea subiendo o sea
accionando el ascensor. Y todas se orientan a entrega o exploración antes que a
la asistencia de personas.

\subsection{Distribución del cómputo}

La última línea trata de dónde debe ejecutarse el procesamiento.
\textcite{kehoe2015cloud} revisan las arquitecturas que lo delegan a servidores
remotos, y \textcite{tahir2025edge} las que lo acercan al borde de la red.
\textcite{malavolta2023offloading} miden el costo real de esa delegación sobre
WiFi en robots terrestres, que es el caso de este proyecto.

Sus resultados pesaron en una decisión concreta. La navegación corre completa a
bordo de cada vehículo, y lo único que se centraliza es la coordinación, que
intercambia mensajes pequeños y corre en un computador portátil. Además, la
interfaz debe funcionar sin salida a internet, porque la red del edificio no
siempre la tiene. La sección~\ref{sec:teo-ros2} explica cómo esa distribución se
refleja en la red de comunicaciones.

\subsection{Aporte de este trabajo frente a los antecedentes}
\label{sec:aporte}

La literatura cubre bien cada componente por separado. Hay flujos de navegación
completos sobre ROS, esquemas de coordinación multirrobot a distintos niveles,
soluciones de movilidad vertical y trabajo sobre la integración del usuario. La
tabla~\ref{tab:antecedentes} resume lo que cada línea aporta a este proyecto y el
supuesto que este proyecto no comparte.

\begin{table}[H]
  \centering\small
  \caption{Líneas de antecedentes, lo que aportan a este trabajo y el supuesto
    que este trabajo cambia.}
  \label{tab:antecedentes}
  \begin{tabular}{@{}p{0.2\textwidth}p{0.36\textwidth}p{0.36\textwidth}@{}}
    \toprule
    \textbf{Línea} & \textbf{Lo que se toma de ella} & \textbf{Supuesto que no se comparte} \\ \midrule
    Robots guía & La integración de navegación e interacción; la localización de Monte Carlo & Un solo robot acompaña todo el recorrido \\ \addlinespace
    Navegación en interiores & El flujo mapa, localización y planificación sobre ROS~2 (Nav2) & Ninguno: se adopta, adaptado a la cinemática del vehículo \\ \addlinespace
    Coordinación multirrobot & La taxonomía de asignación; la planificación con relevos & Lo que se traspasa puede esperar sin límite \\ \addlinespace
    Movilidad entre pisos & El mapa por piso unido por un punto de transición & El robot cruza de nivel \\ \addlinespace
    Distribución del cómputo & La navegación a bordo y la coordinación en un nodo central & --- \\ \bottomrule
  \end{tabular}
\end{table}

Lo que no aparece en ninguna línea es la combinación que aquí se propone:
resolver la discontinuidad entre pisos sin que ningún robot la atraviese. El
recorrido se reparte entre dos agentes dedicados, cada uno a un nivel, y la
persona pasa de uno a otro en la zona de transición.

La diferencia con los trabajos de movilidad vertical es económica antes que
algorítmica. Esos trabajos obtienen la continuidad con hardware: mecanismos de
locomoción para escaleras o intervención del ascensor. Este proyecto la obtiene
con comunicación entre agentes, que es lo que los trabajos de coordinación ya
resuelven, aplicada a un problema donde no se había usado.

% =====================================================================
\section{Marco conceptual}
\label{sec:conceptual}
% =====================================================================
% FORMA: parrafos corridos, no listas de definiciones (pedido del director,
% 2026-10-05): un termino definido dentro de una frase que lo conecta con el
% anterior se entiende mejor que el mismo termino en una ficha.
% Los terminos siguen en negrita para poder barrer la seccion con la vista.
% =====================================================================

Aquí se definen los términos necesarios para leer el resto del documento. El
lector supuesto tiene formación en ingeniería electrónica. Las seis categorías
son las del anteproyecto, y dentro de cada una se añaden los términos que
aparecieron durante la implementación. El fundamento teórico de cada concepto
está en la sección~\ref{sec:teorico}.

\subsection{Navegación y percepción}

Se llama \textbf{navegación autónoma} al desplazamiento de un robot sin
intervención humana. Reúne tres capacidades que trabajan juntas: percibir, saber
dónde se está y decidir hacia dónde moverse. La primera pieza de «saber dónde se
está» es la \textbf{odometría}, que estima el desplazamiento del vehículo desde
su punto de partida a partir de sus propias mediciones. Cuando esas mediciones
salen de comparar barridos sucesivos de un sensor láser, se habla de
\textbf{odometría láser}; es la única disponible en los vehículos de este
proyecto, que no tienen codificadores en las ruedas. Toda odometría se desvía con
el recorrido, así que hay que corregirla contra una referencia externa. De eso
se encarga \textbf{AMCL}, un algoritmo de localización por filtro de partículas
que combina la odometría con las lecturas del láser y sitúa al vehículo dentro de
un mapa conocido \autocite{fox1999monte}.

El sensor láser de barrido se llama \textbf{LiDAR}: mide distancias en un plano,
en 360$^\circ$ alrededor del vehículo. Junto a él, la tarjeta de cada vehículo
incluye una \textbf{unidad de medición inercial} o IMU, con un giroscopio que
mide la velocidad de giro y un acelerómetro. Las dos fuentes se combinan con un
\textbf{filtro de Kalman extendido}, un estimador que pondera cada medición según
su incertidumbre y entrega una sola estimación del movimiento
\autocite{thrun2005probabilistic}.

Hay un concepto que decide si todo lo anterior puede funcionar: la
\textbf{observabilidad} de un estado. Una magnitud es observable cuando se puede
estimar a partir de las mediciones disponibles. La propiedad depende de la
combinación de entorno y sensor, no solo del sensor, de modo que una misma
magnitud puede ser observable en un sitio e inobservable en otro. Cuando es
inobservable, ningún ajuste de parámetros la recupera, porque la información no
está en la señal. Para cuantificarla, el proyecto definió la \textbf{información
de avance} de un barrido: la fracción de los rayos del LiDAR que inciden sobre
superficies orientadas hacia la dirección de marcha. Solo esos rayos cambian de
longitud cuando el vehículo avanza. La sección~\ref{sec:teo-localizacion} desarrolla el
caso.

El entorno se representa de dos formas complementarias. El \textbf{mapa de
ocupación} dice qué celdas están libres y cuáles no. El \textbf{mapa de costos}
asigna a cada celda un costo de tránsito en lugar de una sola etiqueta, de modo
que el planificador puede preferir el centro del pasillo sin que los bordes
queden prohibidos \autocite{lu2014layered}. Hay un mapa de costos global, que
cubre todo el mapa, y uno local, limitado al entorno inmediato del vehículo, que
registra los obstáculos que el láser ve en ese momento.

Sobre esa representación se calcula la trayectoria. En un vehículo que gira con
radio mínimo y puede ir hacia adelante y hacia atrás, las trayectorias más
cortas son las del \textbf{modelo de movimiento Reeds-Shepp}
\autocite{reeds1990optimal}. El punto donde el vehículo invierte el sentido de la
marcha se llama \textbf{cúspide}, y contarlas indica cuántas maniobras hicieron
falta. Al final del recorrido interviene el \textbf{verificador de meta}, que
decide si el vehículo llegó comparando su posición, y a veces su orientación,
contra una tolerancia.

Queda un término que es del edificio y no del vehículo: la \textbf{zona de
transición}, el lugar donde se cambia de nivel. En este sistema es la escalera, y
ningún vehículo la atraviesa.

\subsection{Control y locomoción}

La \textbf{cinemática de robots móviles} es el modelo matemático que relaciona
las velocidades de los actuadores con el movimiento del vehículo
\autocite{siegwart2011introduction}. La plataforma de este trabajo sigue un
\textbf{modelo cinemático tipo automóvil}, también llamado \textbf{dirección
Ackermann}: cuatro ruedas, dirección en el eje delantero y una relación entre el
ángulo del servomotor y el radio de giro. Ese modelo hace que las ruedas interior
y exterior giren en torno al mismo centro instantáneo de rotación, de modo que no
derrapan. Lo que impone, y condiciona todo lo demás, es un \textbf{radio de giro
mínimo}: el vehículo no puede girar sobre su propio eje. Un robot que sí puede
hacerlo, como uno de tracción diferencial, se dice que gira en el sitio.

Esa restricción obliga a elegir con cuidado los dos componentes que calculan y
siguen la trayectoria. La planificación corre a cargo de un \textbf{planificador
híbrido A*}, que busca rutas trabajando sobre posición y orientación a la vez
\autocite{dolgov2010path}. El seguimiento lo hace un \textbf{controlador de
persecución pura regulada}, que persigue un punto móvil sobre la trayectoria y
regula la velocidad según la curvatura y la cercanía de obstáculos
\autocite{macenski2023rpp}.

Las órdenes llegan a los actuadores como señales de \textbf{PWM}, modulación por
ancho de pulso, que controlan el motor de tracción y el servomotor de dirección
variando el ciclo de trabajo. En vehículos con realimentación de velocidad, la
regulación fina se hace con un \textbf{control PID},
proporcional-integral-derivativo. El vehículo de este proyecto no tiene esa
realimentación: su tracción funciona en lazo abierto, y la misma orden produce
velocidades distintas según la carga de la batería. Entre la orden y el
movimiento aparece además la \textbf{banda muerta de tracción}, el rango de
velocidades pedidas en el que el motor no recibe potencia suficiente para mover
el vehículo. En esta plataforma está por debajo de 0,40~m/s.

Por último, la \textbf{media vuelta en varios tiempos} es la maniobra con que un
vehículo de este tipo invierte su rumbo en un pasillo angosto: alterna marcha
atrás y avance con la dirección a tope, hasta completar el giro.

\subsection{Comunicación y sistemas distribuidos}

Un \textbf{middleware robótico} es la capa de software que comunica entre sí los
módulos del sistema. Aquí es \textbf{ROS~2}, que organiza el software en
\textbf{nodos}, programas que cumplen cada uno una función, y los comunica por
\textbf{tópicos} de publicación y suscripción, \textbf{servicios} de petición y
respuesta, y \textbf{acciones} \autocite{macenski2022ros2}. Por debajo, ROS~2
implementa todo eso sobre \textbf{DDS}, un estándar de comunicación en el que los
participantes de la red se descubren solos \autocite{omg2015dds}.

De ese mecanismo de \textbf{descubrimiento} salen varios conceptos que en la
práctica deciden quién habla con quién. El \textbf{dominio} es un número que
segmenta el descubrimiento: dos procesos con dominio distinto no se ven entre
sí. La \textbf{partición de comunicaciones} es un filtro más fino, dentro de un
mismo dominio, que solo deja intercambiar mensajes a los procesos de la misma
partición. El \textbf{espacio de nombres} es un prefijo, como
\texttt{/robot1}, que se antepone a los tópicos, servicios y acciones de un nodo
para distinguir a qué vehículo pertenecen. Y el \textbf{rango de descubrimiento}
limita a quién se anuncia cada proceso: solo a los de su propio computador, o
también a una lista fija de equipos. Con esa última opción la red adopta una
\textbf{topología en estrella}, en la que los vehículos se comunican con el
computador central y no entre sí.

De los patrones de comunicación, el que sostiene el sistema es la
\textbf{acción}, pensada para tareas largas. Tiene una meta, retroalimentación
periódica y un resultado final, y se puede cancelar. El nodo que la ofrece es el
\textbf{servidor de acción} y el que la invoca, el \textbf{cliente}. El
coordinador de este proyecto cumple los dos papeles: es cliente ante cada agente
y servidor ante la interfaz. Para aceptar cancelaciones, un servidor de acción
debe registrar una función de respuesta a la cancelación; si no lo hace, ROS~2
las rechaza todas sin informar del motivo. Cada canal lleva además sus políticas
de entrega, que se llaman \textbf{calidad de servicio} y se configuran por
separado.

La interfaz web no habla ROS, así que entre las dos hay un traductor: el
\textbf{puente \textit{rosbridge}}, un servidor que convierte entre el protocolo
de ROS~2 y mensajes JSON sobre WebSocket \autocite{crick2017rosbridge}. Es lo que
permite que una página web opere el sistema sin instalar nada. Finalmente, un
\textbf{contenedor} es un entorno de ejecución aislado, con su propio sistema de
archivos y sus propias bibliotecas, que corre sobre el sistema operativo del
equipo. Permite tener en un mismo computador dos versiones distintas de ROS~2.

\subsection{Arquitectura del sistema}

La \textbf{asignación de tareas} es el reparto de actividades entre los robots
según su disponibilidad y su ubicación \autocite{gerkey2004taxonomy}. A veces la
tarea no cabe en un solo agente, porque atraviesa la frontera entre dos niveles.
Entonces el reparto toma la forma de un \textbf{protocolo de relevo}: el traspaso
de la responsabilidad de guiar al usuario de un agente a otro en un punto
definido. Ese lugar, que en este edificio es la escalera, se llama \textbf{punto
de transferencia}.

Tres piezas más completan la arquitectura. La \textbf{máquina de estados de
misión} es un conjunto ordenado de estados que describe en qué punto va una
misión; se publican para que la interfaz y el registro de evaluación lo sepan. El
\textbf{planificador puro} calcula los tramos de la misión a partir del catálogo
de puntos y del par origen-destino. Se llama «puro» porque no depende del
middleware ni de ningún estado externo, lo que permite verificarlo con cientos de
casos en milisegundos, sin levantar el simulador. Y el \textbf{registro de
misión} es un documento estructurado que se compone a partir de la grabación de
cada ejecución y se valida contra un esquema. Es la evidencia primaria del
experimento.

\subsection{Interacción humano-robot}

La \textbf{interacción humano-robot}, o HRI por sus siglas en inglés, es el campo
que estudia el diseño y la evaluación de sistemas robóticos pensados para
interactuar con personas \autocite{jimenez2024arquitectura}. Su parte visible
aquí es una \textbf{interfaz gráfica de usuario}, el entorno visual donde la
persona y la máquina se encuentran. El acto que la pone en marcha es la
\textbf{solicitud de servicio}: la petición del usuario para recibir guiado.

Dos decisiones de diseño tienen nombre propio porque se tomaron después de ver el
sistema en uso. La \textbf{interfaz responsiva de selección} permite buscar un
destino escribiendo su nombre, sin perder la opción de ver el catálogo completo.
Con diez destinos la diferencia es pequeña; con cien no lo sería. Y la
\textbf{retroalimentación literal} es un principio: el texto que ve el usuario lo
redacta el componente que conoce el estado del sistema, y la interfaz lo muestra
tal cual. Así la interfaz no puede contradecir al coordinador.

\subsection{Métricas y criterios de evaluación}

Las cuatro métricas del proyecto son las del anteproyecto. El \textbf{tiempo de
respuesta} es lo que pasa entre la solicitud del usuario y el momento en que el
vehículo arranca la tarea. El \textbf{tiempo de asignación} es lo que pasa entre
aceptar la solicitud y designar qué agente la atiende. La \textbf{tasa de éxito}
es la proporción de solicitudes que terminan bien. Y la \textbf{continuidad de
servicio} es la capacidad de mantener la asistencia mientras el usuario cambia de
nivel.

Medir bien exige cuatro conceptos más, que el proyecto incorporó durante la
evaluación. El \textbf{intervalo de confianza de Wilson} estima el intervalo de
una proporción y es más fiable que la aproximación normal cuando la muestra es
pequeña o la proporción se acerca a 0 o a 1 \autocite{brown2001interval}. La
\textbf{falsabilidad} de un criterio es la propiedad de que alguna ejecución del
sistema pueda, al menos en principio, incumplirlo \autocite{popper1959logic}. Un
\textbf{invariante estructural} es una propiedad que se cumple siempre por el
diseño del sistema, de modo que no es un resultado experimental y no debe
presentarse como tal. Y la \textbf{cota superior por misión} resuelve qué hacer
cuando el instrumento no tiene resolución para dar el valor exacto: se informa el
límite que sí se puede garantizar.

La demostración sobre los vehículos reales se organizó en \textbf{pruebas de
verificación}: condiciones comprobables, acordadas con el director antes de
correrlas, que había que superar en orden para pasar a la siguiente etapa. Un
ejemplo es que la odometría registre el avance con un error del 10~\% o menos. Para
decidir si se superaron hace falta una \textbf{verdad de terreno}, la medida
independiente con la que se contrasta el sistema, que en este trabajo es la
distancia tomada con flexómetro. En la campaña con los dos vehículos esa medida
faltó, y la sección~\ref{sec:lim-llegada} explica qué cambia en la lectura de sus
resultados.

% =====================================================================
\section{Marco teórico}
\label{sec:teorico}
% =====================================================================

Cada subsección expone el fundamento de una parte del sistema y la literatura que
lo respalda. El orden va de la coordinación, que es el aporte del trabajo, a la
navegación de cada vehículo, y termina en los criterios con que se evaluó.

\subsection{Asignación de tareas en sistemas multirrobot}
\label{sec:teo-asignacion}

\textcite{gerkey2004taxonomy} clasifican los problemas de asignación de tareas
entre robots según tres ejes. El primero distingue robots que ejecutan una sola
tarea a la vez (ST, \textit{single-task}) de robots que atienden varias (MT). El
segundo distingue tareas que requieren un solo robot (SR, \textit{single-robot})
de tareas que requieren varios (MR). El tercero distingue la asignación
instantánea (IA), que se decide al llegar cada tarea, de la extendida en el
tiempo (TA), que planifica asignaciones futuras.

La asignación de este sistema es ST--SR--IA. Cada vehículo ejecuta un tramo a la
vez, cada tramo lo ejecuta un solo vehículo, y el agente se elige al llegar la
solicitud, sin planificar las siguientes. Los mismos autores demuestran que esa
clase es una instancia del problema de asignación óptima, que el método húngaro
resuelve en tiempo polinómico \autocite{kuhn1955hungarian}. Es decir, el óptimo
está disponible y no es costoso de calcular.

En este problema, sin embargo, la matriz de costos tiene un solo candidato
admisible por tarea. Cada vehículo está confinado a su nivel y hay uno solo
por nivel. Aplicar el método húngaro a esa matriz devuelve siempre la misma
correspondencia que el sistema guarda como tabla de configuración, en la que
cada nivel tiene asignado su vehículo. La ausencia de un optimizador es, por
tanto, una consecuencia de la estructura del problema y no una carencia de la
implementación.

Por la misma razón se descartaron los mecanismos de mercado. En el protocolo de
red de contratos \autocite{smith1980contract}, un gestor anuncia la tarea, los
agentes envían una oferta con su costo y la tarea se adjudica a la mejor.
MURDOCH \autocite{gerkey2002sold} aplica ese esquema a robots sobre un modelo de
publicación y suscripción, y CBBA \autocite{choi2009consensus} lo lleva a una
subasta descentralizada en la que los agentes llegan a consenso sin un
coordinador central. Con un solo agente elegible por tramo, cualquier subasta
tiene un único postor: añade mensajes y latencia sin cambiar la decisión. CBBA,
además, elimina el coordinador central, que la arquitectura mantiene por diseño.

La misión completa sí se sale de esa clase. El relevo impone una restricción de
precedencia: el tramo del segundo vehículo no puede empezar hasta que el primero
termine el suyo y el usuario cambie de piso. Las agendas de los dos agentes
quedan acopladas, y la taxonomía de 2004 está construida para tareas
independientes. \textcite{korsah2013comprehensive} la extienden según el grado
de interdependencia entre tareas, y en su clasificación esta misión pertenece a
la clase XD, de dependencias entre agendas (\textit{cross-schedule
dependencies}). Ahí está el aporte del trabajo: la asignación es trivial y óptima
por construcción, y la coordinación del relevo pertenece a una clase de problema
que la taxonomía clásica del área no contempla.

El análisis también delimita el alcance. El sistema atiende una misión a la vez y
no tiene noción de agente ocupado ni cola de solicitudes. En cuanto haya dos o
más vehículos elegibles por tarea, o misiones simultáneas, la matriz de costos
deja de tener un único candidato y los métodos anteriores pasan a decidir de
verdad. Esa extensión se plantea en los trabajos futuros.

\subsection{ROS~2 y la comunicación entre los componentes}
\label{sec:teo-ros2}

Todos los componentes se comunican mediante ROS~2. \textcite{macenski2022ros2}
describen su arquitectura como un rediseño sobre DDS, el estándar de comunicación
de datos en tiempo real del Object Management Group \autocite{omg2015dds}, y no
como una actualización de su antecesor. DDS aporta el descubrimiento automático
de los participantes y las políticas de calidad de servicio de cada canal.

El descubrimiento automático tiene un costo que crece con el tamaño del sistema.
Por omisión, cada proceso se anuncia en la red y empareja sus publicadores y
suscriptores con los de todos los demás. Con las dos cadenas de navegación
completas encendidas, unos cuarenta procesos se anunciaban unos a otros por la
red inalámbrica, y la carga de la tarjeta de un vehículo pasó de 6 a entre 17 y
33 (sección~\ref{sec:dos-vehiculos}). ROS~2 Jazzy permite restringir el
descubrimiento a los procesos del mismo computador y declarar, además, una lista
fija de equipos con los que sí se comunica. Con esas dos opciones la red adopta
la topología en estrella del diseño: dentro de cada vehículo los procesos se
descubren entre sí, y hacia afuera solo se anuncia al computador portátil lo que
el coordinador necesita. Los vehículos no se ven entre sí.

Dos mecanismos más separan a los vehículos. La partición de comunicaciones deja
dentro de cada vehículo los tópicos del software del fabricante, que no llevan
espacio de nombres. Y un espacio de nombres distinto para cada vehículo
distingue la navegación, el estado y la maniobra de media vuelta de cada agente
cuando llegan al portátil.

El sistema usa además dos distribuciones de ROS~2. La simulación corre sobre
Humble, porque los paquetes que integran ROS~2 con la versión de Gazebo empleada
no están publicados para Jazzy. Los vehículos corren sobre
Jazzy, la que corresponde al sistema operativo de sus tarjetas, Ubuntu 24.04. El
código fuente es el mismo en los dos destinos. Lo que cambió entre
distribuciones es la definición del mensaje de la acción de navegación, de modo
que un coordinador compilado en Humble no puede enviar metas a un agente
compilado en Jazzy. Por eso
el coordinador corre siempre en la distribución de los agentes que comanda: en
simulación, sobre Humble; con los vehículos, sobre Jazzy, dentro de un contenedor
en el portátil. La configuración de Nav2 tiene una variante por distribución,
derivada de una base común.

\subsection{Cinemática tipo automóvil y planificación de trayectorias}
\label{sec:teo-planificacion}

Un vehículo con dirección Ackermann está sujeto a una restricción no holonómica:
no puede desplazarse de lado ni girar sobre su propio eje, y su curvatura está
acotada por el radio de giro mínimo \autocite{siegwart2011introduction}. La
consecuencia práctica es que no toda trayectoria libre de obstáculos es
ejecutable. Un planificador que solo considera la posición propone giros que el
vehículo no puede hacer.

\textcite{reeds1990optimal} resolvieron el problema geométrico de fondo: el
camino más corto entre dos poses para un automóvil que puede avanzar y
retroceder, con radio de giro mínimo. Demostraron que basta considerar caminos
formados por arcos de radio mínimo y segmentos rectos, con a lo sumo dos
inversiones de marcha, y que el óptimo está siempre entre un conjunto finito de
candidatos calculables con fórmulas explícitas. Como el vehículo de este trabajo
puede retroceder, ese es su modelo de movimiento.

Sobre un mapa con obstáculos hace falta además una búsqueda.
\textcite{dolgov2010path} propusieron el A* híbrido: una búsqueda A* sobre el
estado continuo del vehículo (posición y orientación), en la que cada expansión
respeta la cinemática y el camino resultante se suaviza después. Lo validaron en
los estacionamientos del DARPA Urban Challenge de 2007, un escenario de
maniobras con inversión de marcha parecido al de un pasillo angosto. Nav2 lo
implementa en el planificador Smac, que admite el modelo Reeds-Shepp y está
pensado para robots tipo automóvil \autocite{macenski2024smac}. Es el planificador
que usan los vehículos de este proyecto.

El planificador trabaja sobre mapas de costos. \textcite{lu2014layered}
propusieron organizarlos por capas, cada una con una fuente de información. La
capa estática viene del mapa y la de obstáculos, de las lecturas del sensor. La
de inflado extiende el costo alrededor de cada obstáculo según el tamaño del
vehículo. El mapa de costos global combina las tres y sirve al planificador; el
local cubre los alrededores del vehículo y sirve al controlador. La configuración
de este trabajo usa esas tres capas.

\subsection{Seguimiento de la trayectoria}
\label{sec:teo-seguimiento}

Una vez calculada la trayectoria, el controlador la convierte en órdenes de
velocidad. La persecución pura \autocite{coulter1992implementation} es un
controlador geométrico. Elige sobre la trayectoria un punto situado a una
distancia de anticipación $\ell$ del vehículo, y calcula el arco de circunferencia
que lleva al vehículo hasta él. Llamando $y$ a la distancia lateral del punto al
eje del vehículo, la curvatura ordenada es
\begin{equation}
  \kappa = \frac{2\,y}{\ell^{2}} .
\end{equation}
Es simple, no requiere un modelo dinámico y se comporta bien en vehículos tipo
automóvil, porque siempre ordena arcos que la dirección puede seguir.

\textcite{macenski2023rpp} le añadieron una regulación de velocidad pensada para
robots de servicio. El controlador reduce la velocidad en las curvas cerradas y
cerca de los obstáculos. Además admite la marcha atrás, de modo que sigue las
cúspides que propone el planificador. En este trabajo se configuró con una
velocidad deseada de 0,5~m/s, una distancia de anticipación entre 0,3 y 0,9~m,
marcha atrás habilitada y el giro en el sitio deshabilitado, porque el vehículo no
puede hacerlo.

El desarrollo permitió caracterizar un límite de esa configuración. Cuando una
cúspide queda muy cerca de la meta, la distancia de anticipación se acorta, y por
debajo de cierto umbral geométrico una guarda de seguridad del propio controlador
anula la curvatura calculada. El vehículo sigue invirtiendo la marcha, pero ya no
dirige, y maniobra decenas de veces en el mismo punto. El origen no está en el
planificador ni en la coordinación, sino en un controlador de referencia que
trabaja en el borde de su dominio de validez. Por eso la media vuelta se
resolvió con una maniobra propia del agente (sección~\ref{sec:media-vuelta}).

La banda muerta de tracción agrega otro límite. Como el vehículo no se mueve por
debajo de 0,40~m/s, la velocidad mínima de aproximación del controlador se fijó
en ese valor, y el vehículo llega a la meta sin una fase de aproximación lenta
(sección~\ref{sec:banda-muerta}).

\subsection{Localización, odometría y observabilidad}
\label{sec:teo-localizacion}

La localización de Monte Carlo representa la creencia sobre la pose del vehículo
con un conjunto de partículas, cada una con un peso
\autocite{fox1999monte,thrun2005probabilistic}. En cada paso, el modelo de
movimiento desplaza las partículas según la odometría. Después, el modelo de
medición las pondera según cuánto se parece el barrido del láser a lo que se
vería desde esa pose en el mapa. Por último, un remuestreo concentra las
partículas en las poses más probables. AMCL es la variante que ajusta el número de partículas a la
incertidumbre; en este trabajo, entre 500 y 2000. Su precisión depende de la
calidad de la odometría entre una actualización y la siguiente.

Como los vehículos no tienen codificadores en las ruedas, la odometría sale del
láser. \textcite{jaimez2016planar} estiman el movimiento plano de un láser a
partir de barridos consecutivos con un enfoque de flujo de distancias (\textit{range
flow}). Alinean los barridos de forma densa, usando sus gradientes, sin buscar
correspondencias entre puntos, y tardan alrededor de un milisegundo por barrido. Es el método que implementa el paquete rf2o, que usan los vehículos.

Todo estimador que se apoya en la geometría del entorno hereda una limitación.
\textcite{zhang2016degeneracy} la formalizan como degeneración: hay direcciones
del espacio de estados en las que las mediciones no restringen la solución, y en
ellas el estimador no tiene información. Un pasillo recto y de sección uniforme
es el caso típico. Al avanzar a lo largo de su eje, cada barrido es casi idéntico
al anterior, y el estimador devuelve un desplazamiento cercano a cero. El
resultado es coherente con lo que mide el sensor, porque desde él no ha cambiado
nada. La posición longitudinal es inobservable, y ningún ajuste de parámetros la
recupera.

De ahí sale la diferencia entre dos tipos de error que se parecen. Un error de
calibración se corrige ajustando parámetros. Una inobservabilidad estructural no
se corrige así, porque la información no está en la señal. La información de
avance, definida en la sección~\ref{sec:conceptual}, mide cuánta hay en cada sitio. El
criterio que el proyecto adoptó para elegir dónde operar es que el entorno ofrezca
estructura no uniforme dentro del alcance del sensor durante todo el recorrido.

La construcción de mapas comparte el mismo límite. La herramienta de SLAM de
ROS~2 \autocite{macenski2021slam} construye el mapa a partir de la misma
odometría del láser, y en un pasillo uniforme no tiene con qué cerrarlo. Por eso
los mapas de este proyecto se derivaron de un levantamiento con flexómetro
(sección~\ref{sec:mapas}).

El rumbo tiene una solución distinta, porque el giroscopio de la IMU lo mide de
forma directa, sin depender del entorno. Un filtro de Kalman extendido combina
las dos fuentes \autocite{thrun2005probabilistic}. El paquete robot\_localization
lo implementa para ROS de forma general, permitiendo elegir qué componente de cada
sensor se fusiona \autocite{moore2014generalized}. En este trabajo el filtro toma
de la odometría láser las velocidades de avance y de la IMU la velocidad de giro.
El rumbo mejora; la posición no, porque sigue saliendo del láser.

\subsection{Simulación del sistema}
\label{sec:teo-simulacion}

La campaña de evaluación corrió en Gazebo, un simulador de código abierto pensado
desde su origen para varios robots, con motor de física, modelos de sensores e
integración con ROS \autocite{koenig2004gazebo}. Se usó la versión 11 de Gazebo,
integrada con ROS~2 Humble. La cinemática Ackermann del vehículo se simula con un
complemento propio del paquete de simulación del DeepRacer, en lugar del genérico
de Gazebo.

La simulación levanta un proceso de simulación por vehículo y un solo dominio de
comunicaciones para los dos. Los agentes se separan por espacio de nombres,
incluido el marco de referencia de su mapa. Dar a cada vehículo un dominio
distinto parece más simple, pero con la configuración por omisión corta el
descubrimiento entre el coordinador y los agentes: separar por dominios y
comunicar a los agentes con el coordinador resultan mutuamente excluyentes.

El tiempo que ven los nodos durante la simulación avanza en saltos de 100~ms, y
eso fija la resolución con que se pueden medir los tiempos de la campaña. El efecto sobre las métricas se
discute en la sección~\ref{sec:limitaciones}.

\subsection{La interfaz web como cliente del sistema}
\label{sec:teo-interfaz}

\textcite{crick2017rosbridge} propusieron \textit{rosbridge} para que programas
ajenos a ROS, como un navegador, puedan usar tópicos, servicios y acciones a
través de mensajes JSON sobre WebSocket. La versión para ROS~2 del puente expone
las acciones con operaciones propias para enviar la meta, recibir la
retroalimentación y recibir el resultado. La biblioteca de cliente más difundida
del ecosistema web de ROS implementa las acciones de la primera generación de
ROS, que no son compatibles con esas operaciones. Por eso la interfaz de este
trabajo es un cliente propio, escrito sobre el protocolo del puente. Tampoco usa
dependencias externas, porque debe funcionar sin salida a internet.

Un segundo principio rige la interfaz: no contiene lógica del dominio. El texto
que lee el usuario en cada estado lo escribe el coordinador, y la interfaz lo
muestra literal. Es la separación de responsabilidades entre el dominio y la
presentación. Si el navegador interpretara el estado para decidir qué texto
mostrar, se duplicaría en el cliente una lógica que ya vive en el planificador, y
las dos copias acabarían divergiendo.

\subsection{Evaluación experimental}
\label{sec:teo-evaluacion}

La evaluación se apoya en las cuatro métricas del anteproyecto, cada una con una
definición operativa que indica qué marca del registro se resta de cuál.
El protocolo se escribió antes de instrumentar el sistema, y las misiones se
sortean con una semilla fijada antes de correrlas. Las dos decisiones tienen el
mismo propósito: que ni el instrumento ni los recorridos se elijan después de ver
los resultados.

La tasa de éxito es una proporción estimada sobre una muestra pequeña, y para
ella el intervalo de confianza habitual, el de la aproximación normal, falla.
\textcite{brown2001interval} mostraron que su cobertura real oscila de forma
irregular, incluso con muestras grandes, y empeora con proporciones cercanas a 0
o a 1. Recomiendan el intervalo de
\textcite{wilson1927probable} para muestras pequeñas. Sea $n$ el número de
misiones, $\hat{p}$ la proporción observada de éxitos y $z$ el cuantil de la
confianza deseada ($z = 1{,}96$ para el 95~\%). El intervalo de Wilson es
\begin{equation}
  \frac{\hat{p} + \dfrac{z^{2}}{2n} \pm z\sqrt{\dfrac{\hat{p}(1-\hat{p})}{n} + \dfrac{z^{2}}{4n^{2}}}}{1 + \dfrac{z^{2}}{n}} .
\end{equation}
Con las treinta misiones de la campaña y veintiséis éxitos, da entre 70,3~\% y
94,7~\%. El ancho del intervalo decide qué se puede afirmar: con esa muestra se
sostiene que la tasa de éxito supera el 70~\%, no que valga 86,7~\%.

El segundo principio es de método. Para \textcite{popper1959logic}, una
afirmación aporta conocimiento empírico solo si alguna observación posible podría
refutarla. Trasladado a un experimento, un criterio de aceptación solo es
evidencia si alguna ejecución del sistema puede incumplirlo. En este trabajo hubo
un caso concreto. El criterio de la variable de respuesta principal, la
continuidad del servicio, estaba definido de forma que ninguna ejecución podía
incumplirlo: los únicos puntos del código que producen un incumplimiento ocurren
fuera de la ventana en que el criterio se evalúa. Treinta
verificaciones satisfactorias de un criterio así no muestran que el sistema
cumpla la propiedad. La corrección no fue inventar una ruta de fallo, que habría
sido instrumentar el experimento para que pareciera más informativo. Se
reformuló la afirmación como un invariante estructural del diseño y se respaldó
con una magnitud que sí varía: el intervalo de tiempo entre el final del tramo
de un vehículo y el inicio del otro (sección~\ref{sec:limitaciones}).

La demostración sobre los vehículos reales sigue otra lógica. La norma ASTM
F3244-21 \autocite{astm2021f3244} establece un método de prueba para la
navegación de robots en un área definida, con repeticiones del recorrido entre
marcas fijas en el piso. El número de repeticiones fija
cuánto se puede afirmar. Si las $n$ repeticiones salen bien, la fiabilidad mínima
que se demuestra con una confianza $C$ es
\begin{equation}
  R = (1-C)^{1/n} ,
\end{equation}
que con $n = 10$ y $C = 0{,}85$ da $R \approx 0{,}83$. Con las cinco a diez
repeticiones que caben en el calendario, la campaña física demuestra que el
sistema opera sobre los vehículos, pero no estima su tasa de éxito con la
precisión de la simulación. Por eso las dos campañas se reportan por separado:
la de simulación produce la evidencia estadística y la física, la demostración. En
los vehículos el protocolo mide la llegada con flexómetro y la tolerancia es de 0,5~m,
fijada por el director para las pruebas con los vehículos; la simulación
conserva 0,25~m.

% =====================================================================
\section{Marco legal}
\label{sec:legal}
% =====================================================================

El proyecto se desarrolla dentro de la política pública colombiana de
transformación digital y de ciencia, tecnología e innovación, que enmarcan dos
documentos. El CONPES 3975 \autocite{conpes3975} fija la Política Nacional para la
Transformación Digital e Inteligencia Artificial. El CONPES 4069
\autocite{conpes4069} establece la Política Nacional de Ciencia, Tecnología e
Innovación. Ambos piden a la academia que desarrolle prototipos con herramientas
de la Cuarta Revolución Industrial, entre ellas la robótica autónoma, y que lo
haga con diseño local e impacto social.

Este trabajo encaja en esa línea. Es un prototipo de asistencia robótica
construido sobre plataformas de bajo costo y software abierto, dentro de las
limitaciones presupuestales de un entorno universitario.

La presentación del propio documento sigue la norma NTC 1486 del ICONTEC
\autocite{ntc1486}, que regula cómo se presentan las tesis y los trabajos de grado y
que la plantilla institucional declara como su base.

\subsection{Alcance de la operación y responsabilidad}

El alcance del proyecto excluye de forma explícita que los robots interactúen
mecánicamente con el edificio. No accionan puertas ni botones de ascensor, no
transportan carga y no cruzan las zonas de transición vertical. Esa delimitación
se adoptó por razones técnicas y de costo, y tiene además una consecuencia
normativa favorable: evita intervenir sistemas sujetos a regulación específica de
seguridad en edificaciones.

Las pruebas se hicieron en pasillos de uso común, en franjas de baja circulación,
con una persona junto a cada vehículo durante las maniobras y la posibilidad de
cancelar la misión desde la interfaz en cualquier momento.

% =====================================================================
\section{Marco ambiental}
\label{sec:ambiental}
% =====================================================================

El impacto ambiental del prototipo se concentra en tres frentes: el hardware que
emplea, la energía que consume y el destino de sus componentes al final de la
vida útil.

\subsection{Reutilización de hardware existente}

Las dos plataformas robóticas son equipo institucional que ya existía. El
proyecto no las compró, y el presupuesto del anteproyecto las contabiliza como
uso y desgaste. La decisión se tomó por razones económicas, pero reduce también
la huella material del trabajo: no entran plataformas nuevas al inventario y se
alarga la vida útil de equipos que estaban subutilizados.

Lo mismo ocurre con el entorno de evaluación. Es el propio edificio de la
universidad, sin construir infraestructura dedicada ni modificar la existente. El
sistema opera sobre la planta tal como está.

\subsection{Consumo energético y autonomía}

Cada vehículo lleva dos fuentes de alimentación independientes, una para el
cómputo y otra para la tracción. Un tramo de navegación dura entre unos segundos
y un minuto, y una sesión de campo completa gasta una fracción de la carga
disponible.

El consumo de los vehículos durante toda la campaña es pequeño frente al de la
estación de trabajo que ejecuta la simulación, que corre dos instancias del
simulador a la vez y de forma sostenida. En un proyecto de robótica de esta
escala, el mayor consumo energético está en la simulación y no en los robots. Ese
hecho respalda la estrategia de verificación que se adoptó, en la que la lógica de
coordinación se comprueba como función pura, sin simulador, en milisegundos y con
un consumo despreciable.

\subsection{Ciclo de vida de las baterías}

Las baterías de polímero de litio que alimentan la tracción son el componente más
sensible del sistema desde el punto de vista ambiental. Contienen litio y se
degradan con los ciclos de carga. Al final de su vida útil deben ir a los canales
de disposición de residuos electrónicos de la institución, nunca a la basura
doméstica.

% VERIFICAR: si la universidad tiene un programa institucional de gestión de RAEE
% (residuos de aparatos eléctricos y electrónicos), nombrarlo aquí en concreto.
% No se afirma que exista sin comprobarlo.

El nivel de carga influye además en el comportamiento del vehículo. Como la
tracción funciona en lazo abierto, una misma orden produjo velocidades distintas
en corridas consecutivas, y la batería es la causa más probable. Esa dependencia
obliga a revisar el nivel de carga entre mediciones. Es una buena práctica
experimental que además desincentiva el uso de baterías degradadas.
